% Versión en español. Formato de dos columnas estilo SAI/IEEE (IEEEtran.cls no
% está instalado localmente, así que emulamos la plantilla SAI_PAPER_FORMAT en
% la clase article: dos columnas, tamaño carta, encabezados con numeración
% romana, "Resumen--"/"Palabras clave--" en línea y referencias numéricas en
% orden de cita).
\documentclass[10pt,twocolumn,letterpaper]{article}

\usepackage[utf8]{inputenc}
\usepackage[T1]{fontenc}
\usepackage{lmodern} % fuentes escalables, requeridas por la expansión de microtype
\usepackage[spanish,es-noshorthands]{babel}
\usepackage[letterpaper,top=0.75in,bottom=1in,left=0.625in,right=0.625in,columnsep=0.25in]{geometry}
\usepackage{amsmath,amssymb}
\usepackage{graphicx}
\usepackage{booktabs}
\usepackage{multirow}
\usepackage{caption}
\usepackage{microtype}
\usepackage{xcolor}
\usepackage[numbers,sort&compress]{natbib}
\usepackage{xurl} % permite cortar URLs largas de las referencias
\usepackage[hidelinks]{hyperref}

% Numeración estilo IEEE/SAI: secciones romanas (I, II, ...), subsecciones con letra.
\renewcommand{\thesection}{\Roman{section}}
\renewcommand{\thesubsection}{\Alph{subsection}}
\renewcommand{\thesubsubsection}{\arabic{subsubsection}}
\renewcommand{\tablename}{Tabla}
\captionsetup{font=small,labelsep=period}

\newcommand{\gcp}{\textsc{gcp}}
\newcommand{\ata}{\textsc{a2a}}
\newcommand{\bigoh}{\mathcal{O}}

\title{\textbf{La Federación de Contexto con Lectura Primero y Control por Rol\\
Supera al Paso de Mensajes Punto a Punto en\\
Sistemas Multiagente Basados en LLM:\\
Una Evaluación Controlada por Equidad}}

\author{%
Adrian Issac Mamani Quispe\\
\textit{Universidad Nacional de San Agust\'in de Arequipa}\\
Arequipa, Per\'u\\
admamaniq@unsa.edu.pe
\and
Ryan Fabian Valdivia Segovia\\
\textit{Universidad Nacional de San Agust\'in de Arequipa}\\
Arequipa, Per\'u\\
rvaldiviase@unsa.edu.pe
}

\date{Junio de 2026}

\begin{document}
\maketitle

\begin{abstract}
Los sistemas multiagente construidos sobre grandes modelos de lenguaje (LLM)
coordinan, en su inmensa mayoría, mediante \emph{paso de mensajes}: cada agente
envía la tarea y el contexto directamente a pares con nombre, como en el
protocolo Agent2Agent (A2A) de Google y en los lenguajes clásicos de
comunicación entre agentes de los que desciende. Sostenemos que un sustrato
alternativo ---\emph{federación de contexto con lectura primero y control por
rol} (el Graph Context Protocol, \gcp{})--- es estructuralmente superior en tres
ejes con igual éxito de tarea: \textbf{acoplamiento}, \textbf{fuga de
información} y \textbf{costo}. En \gcp{}, un agente no empuja contexto hacia sus
pares; lo \emph{lee} desde nodos de conocimiento de propiedad independiente a
través de una compuerta de política que autoriza la lectura por rol y registra
una entrada de procedencia por cada decisión. Hacemos la comparación
\emph{justa} manteniendo todo constante salvo la capa de interoperabilidad: ambos
brazos ejecutan el mismo cerebro de agente, el mismo modelo y la misma tarea;
solo difieren la fábrica de herramientas por par y el alojamiento de nodos (un
manejador de \texttt{fetch} de \gcp{} en proceso frente a servidores A2A reales
en \texttt{localhost}). En tres escenarios ---un mercado (acoplamiento), una
organización de software (fuga) y una cadena de suministro entre propietarios
(entre dueños)--- encontramos: (i) el número de aristas de integración por pares
crece como $\bigoh(n^2)$ para \ata{} frente a $\bigoh(n)$ para \gcp{} ($2450$
frente a $50$ en $n{=}50$), con esfuerzo de integración marginal constante para
\gcp{}; (ii) la completitud de procedencia es $1{,}0$ para \gcp{} y $0{,}0$ para
\ata{} en todos los escenarios y semillas, porque las decisiones de lectura se
auditan por construcción mientras que los mensajes punto a punto no dejan rastro
de acceso; y (iii) \gcp{} nunca cuesta más y, en la tarea entre propietarios,
cuesta menos tokens con mucha menor varianza ($707{,}6{\pm}3{,}5$ frente a
$901{,}4{\pm}229{,}2$). El éxito de tarea es igual en ambos brazos. Posicionamos
\gcp{} frente a cinco décadas de teoría del acoplamiento e investigación en
control de acceso, y frente a trabajo reciente sobre seguridad de agentes LLM, y
liberamos el banco de pruebas, los escenarios y las semillas. La evaluación es
preliminar ---un modelo de capa gratuita, cinco semillas--- y somos explícitos
sobre qué autoriza y qué no.
\end{abstract}

\smallskip
\noindent\textbf{\textit{Palabras clave---}}\textit{federación de contexto;
sistemas multiagente; grandes modelos de lenguaje; control de acceso; mínimo
privilegio; procedencia; acoplamiento de software; protocolos de comunicación
entre agentes}
\smallskip

\section{Introducción}

El patrón de coordinación dominante para los sistemas multiagente basados en LLM
es el paso directo de mensajes. Marcos como AutoGen~\citep{wu2023autogen},
ChatDev~\citep{qian2024chatdev} y MetaGPT~\citep{hong2024metagpt} componen
aplicaciones a partir de agentes que envían y reciben mensajes, y el emergente
protocolo Agent2Agent (A2A)~\citep{a2a2025} estandariza exactamente esto: un
agente cliente anuncia y delega tareas a agentes remotos con nombre sobre
JSON-RPC. Este paradigma desciende directamente de los lenguajes clásicos de
comunicación entre agentes KQML~\citep{finin1994} y FIPA-ACL~\citep{fipa2002},
que ya asumían conexiones por pares y protocolos compartidos, prenegociados,
entre cada par que interactúa.

El paso de mensajes tiene un costo estructural que la literatura de ingeniería de
software comprende desde hace cincuenta años. Conectar $n$ componentes por pares
requiere del orden de $n^2$ enlaces; introducir un intermediario compartido al
que cada componente se conecta una sola vez reduce esto a $n$ enlaces. Este es el
``espagueti de integración'' que el catálogo de Patrones de Integración
Empresarial~\citep{hohpe2003} aborda con corredores de mensajes y concentradores,
y es el mismo fenómeno que la teoría del
acoplamiento~\citep{parnas1972,stevens1974,briand1999} formaliza: la
modificabilidad de un sistema está gobernada por el número y la fuerza de sus
dependencias entre componentes.

Estudiamos un sustrato alternativo para la coordinación multiagente, el Graph
Context Protocol (\gcp{}). \gcp{} es de \emph{lectura primero}: un agente obtiene
lo que necesita leyendo contexto desde nodos de conocimiento de propiedad
independiente, en lugar de recibir mensajes empujados. Tiene \emph{control por
rol}: cada lectura pasa por una compuerta de política que autoriza la solicitud
según el rol del solicitante y sus capacidades requeridas, denegando por
defecto~\citep{saltzer1975,sandhu1996}. Y es \emph{auditado}: cada decisión de
lectura ---concesión o denegación--- produce un registro de procedencia
estructurado que nombra al principal, el nodo objetivo, la decisión y la política
coincidente.

\paragraph{Tesis.} Con igual éxito de tarea, la federación de contexto con
lectura primero y control por rol supera al paso de mensajes punto a punto en
acoplamiento, fuga y costo.

\paragraph{Contribuciones.}
\begin{enumerate}
  \item \textbf{Arquitectura.} Describimos \gcp{}: federación de contexto con
  lectura primero y control por rol, con procedencia por lectura de primera clase
  (Sección~\ref{sec:arch}), y la contrastamos con el paso de mensajes basado en
  empuje.
  \item \textbf{Un método de evaluación controlado por equidad.} El defecto más
  común en las comparaciones de arquitecturas es que los brazos difieren en más
  que la arquitectura. Nuestro banco de pruebas construye ambos brazos a partir de
  una sola definición de escenario y un solo modelo inyectado, intercambiando
  \emph{únicamente} la capa de interoperabilidad (Sección~\ref{sec:method}).
  \item \textbf{Un enfrentamiento empírico} sobre tres escenarios con un LLM real
  y cinco semillas, reportando conteos (conexiones, mensajes, tokens) y no solo
  tiempo de reloj (Sección~\ref{sec:results}).
  \item \textbf{Definiciones de métricas} para el acoplamiento estructural, la
  fuga por coincidencia exacta de canario (una cota inferior) y la completitud de
  procedencia (Sección~\ref{sec:metrics}).
\end{enumerate}

Somos deliberadamente conservadores con las afirmaciones empíricas: los
resultados de comportamiento provienen de un único modelo de capa gratuita sobre
cinco semillas y deben leerse como un \emph{piloto}. La
Sección~\ref{sec:threats} expone las amenazas a la validez en su totalidad,
incluido un caso en el que la fuga arquitectónica a favor de \gcp{} es
\emph{latente y no observada} porque el modelo nunca ejercitó la ruta de fuga.

\section{Antecedentes y Trabajo Relacionado}
\label{sec:related}

\subsection{El paso de mensajes como línea base multiagente}
Los lenguajes de comunicación entre agentes establecieron el paradigma del paso
de mensajes: KQML~\citep{finin1994} envolvía cargas de conocimiento en
performativos tipados intercambiados sobre conexiones por pares, y
FIPA-ACL~\citep{fipa2002} estandarizó un mensaje de 12 campos en el que el emisor
y el receptor deben compartir lenguaje de contenido y ontología antes de que
cualquier intercambio tenga éxito. Los marcos modernos de LLM reinventaron esto
en lenguaje natural. AutoGen~\citep{wu2023autogen} programa aplicaciones como
conversaciones entre agentes; ChatDev~\citep{qian2024chatdev} dispone pares de
roles en una ``cadena de chat'' fija; y los estudios del
campo~\citep{guo2024survey} taxonomizan la estructura de comunicación en
variantes punto a punto, centralizada, por capas y de fondo compartido. El Model
Context Protocol~\citep{mcp2024} estandariza el acceso a contexto de
agente a herramienta (una forma temprana de lectura primero), mientras que
A2A~\citep{a2a2025} estandariza la delegación de tareas de agente a agente ---la
línea base con la que comparamos. Trabajo reciente ha comenzado a cuantificar la
sobrecarga del paso de mensajes: \citet{zhang2024agentprune} reportan un exceso
de tokens del 28--73\% en los grafos de comunicación entre agentes y recuperan la
mayor parte podando la topología de mensajes. De manera muy relacionada,
\citet{mao2026delm} reemplazan tanto la orquestación centralizada como el
enrutamiento punto a punto por un contexto compartido verificado y reportan
resultados de vanguardia con aproximadamente la mitad del costo ---evidencia de
que un sustrato de lectura desde contexto compartido puede dominar al paso de
mensajes. A esa línea añadimos autorización por rol, procedencia auditada por
lectura y un enfrentamiento controlado por equidad.

\subsection{Acoplamiento y el problema $n^2$}
Que menos enlaces entre componentes producen sistemas más mantenibles es de los
resultados más antiguos de la ingeniería de software. \citet{parnas1972} mostró
que la descomposición por ocultamiento de información ---exponer solo una interfaz
mínima--- produce un acoplamiento drásticamente menor que la partición por
diagrama de flujo. \citet{stevens1974} introdujeron el acoplamiento y la cohesión
como propiedades medibles y ordinales de la dependencia entre módulos.
\citet{briand1999} dieron un marco teórico de medición y de teoría de grafos en
el que el acoplamiento es una función monótona del conteo de aristas de
dependencia, por lo que una topología de $\bigoh(n)$ aristas está demostrablemente
menos acoplada que una de $\bigoh(n^2)$. \citet{hohpe2003} lo trasladaron a la
integración distribuida: una malla punto a punto completa de $n$ sistemas necesita
$n(n{-}1)/2$ canales, que un corredor de tipo concentrador reduce a $n$. La
investigación en microservicios~\citep{panichella2021} lleva el mismo argumento de
conteo de aristas a los sistemas nativos de la nube, y en el ámbito de los LLM
\citet{shen2025topology} muestran empíricamente que las topologías de agentes
densas y completamente conectadas amplifican la propagación de errores mientras
que las topologías más dispersas la suprimen. Aplicamos esto directamente: \ata{}
crea aristas agente$\leftrightarrow$agente por pares; \gcp{} crea aristas
agente$\rightarrow$almacén.

\subsection{Mínimo privilegio, RBAC y contexto con compuerta}
La compuerta de lectura de \gcp{} operacionaliza el principio de mínimo
privilegio, articulado por \citet{saltzer1975} y equiparado allí con la necesidad
de conocer (need-to-know) militar. El control de acceso basado en
roles~\citep{sandhu1996} desacopla los permisos de los individuos asignándolos a
roles, y el control de acceso basado en atributos~\citep{nist800162} generaliza la
política a atributos arbitrarios; la protección basada en
capacidades~\citep{dennis1966} liga la autoridad a fichas infalsificables y
explícitamente pasadas ---el modelo que siguen las credenciales de rol de \gcp{}.
Las obligaciones de minimización de datos y limitación de
finalidad~\citep{gdpr2016} convierten ``leer solo lo que tu rol requiere'' en una
restricción legal, no en una mera preferencia de diseño. La línea base punto a
punto no tiene una compuerta análoga: la prevención de fugas recae enteramente en
la discreción del emisor.

\subsection{Fuga, procedencia y gobernanza en sistemas de agentes}
La literatura de seguridad establece que la confidencialidad en los sistemas de
agentes es una propiedad del sistema, no solo del modelo~\citep{evertz2024whispers},
que los canales de mensajes entre agentes son el vector de fuga dominante
---filtrando mucho más que las salidas finales y siendo en gran medida invisibles a
la auditoría solo de salidas~\citep{elyagoubi2026agentleak}--- y que restringir lo
que un agente puede \emph{ver} ---en lugar de filtrar lo que puede \emph{decir}---
es la intervención más robusta~\citep{bagdasarian2024airgap}. El Top 10 de OWASP
para aplicaciones LLM~\citep{owasp2025} nombra la Divulgación de Información
Sensible y la Agencia Excesiva como riesgos primarios, con ``adoptar el mínimo
privilegio'' como mitigación normativa. Del lado de la rendición de cuentas, el
control de privilegios determinista en la frontera de herramientas~\citep{shi2025progent},
la delegación de autoridad autenticada y auditable a los
agentes~\citep{south2025delegation} y los registros de auditoría para
LLM~\citep{ojewale2025audit} abogan por registros a prueba de manipulación y
consultables de lo que hizo un sistema. La contribución de \gcp{} a esta línea es
hacer de la \emph{decisión por lectura} el primitivo de procedencia: cada lectura
de contexto autorizada o denegada se registra, algo que el paso de mensajes no
hace.

\paragraph{Diferencia (delta).} El trabajo previo aporta las piezas ---acceso de
lectura primero (MCP), control por rol/atributo (RBAC/ABAC, RAG con compuerta) y
registros de auditoría--- pero no su conjunción, y no un enfrentamiento controlado
por equidad contra una línea base real de paso de mensajes. \gcp{} combina acceso
de lectura primero, control por rol y procedencia auditada por lectura, y medimos
la combinación contra \ata{}.

\section{El Graph Context Protocol}
\label{sec:arch}

\gcp{} modela un ecosistema multiagente como \emph{nodos de conocimiento} y
\emph{nodos de agente} de propiedad independiente. La coordinación es de lectura
primero: un agente adquiere contexto emitiendo una \emph{consulta de contexto} a
un nodo par, que devuelve el contexto que el solicitante está autorizado a leer.
Tres propiedades distinguen al sustrato.

\paragraph{Lectura primero.} Los productores no empujan hacia los consumidores. Un
nodo de conocimiento expone contexto; los consumidores lo extraen bajo demanda.
Ningún agente necesita conocer la identidad de los \emph{consumidores} de sus
productores, y añadir un consumidor no requiere ningún cambio en el productor.
Estructuralmente, cada agente se conecta a los nodos que lee, no a todos los demás
agentes.

\paragraph{Control por rol.} Cada nodo de conocimiento porta una política de
acceso (\texttt{readableByRoles}, \texttt{requiredCapabilities},
\texttt{fallbackAllowed}, \texttt{denialMode}) y el servidor autoriza cada
consulta contra el principal del solicitante, denegando por defecto cuando ninguna
política concede el acceso. La autorización ocurre \emph{antes} de que cualquier
contenido llegue a la ventana de contexto del agente ---el análogo de lectura
primero de la mediación completa~\citep{saltzer1975}.

\paragraph{Auditado.} Cada decisión de lectura emite un registro de procedencia
estructurado (\texttt{principalId}, \texttt{targetNodeId}, \texttt{timestamp},
\texttt{decision}, \texttt{matchedPolicy}) a un sumidero de auditoría enchufable.
Esto produce un libro mayor completo de quién-leyó-qué-cuándo para la federación,
incluidas las denegaciones.

La línea base de paso de mensajes (\ata{}) implementa las mismas tareas de
escenario haciendo que los agentes mantengan referencias a pares mediante tarjetas
e intercambien contexto a través de mensajes directos, con permitir/denegar grueso
declarado en la tarjeta y sin sumidero de auditoría por lectura ---un control fiel
y de buena fe, construido con el mismo cuidado que el brazo \gcp{}.

\section{Modelo de Amenazas y Costo, y Definición de Métricas}
\label{sec:metrics}

Reportamos métricas en tres familias.

\paragraph{Acoplamiento estructural (determinista).} Para $n$ nodos pares, las
\emph{conexiones por pares} son el número de aristas de integración que la
topología requiere. El brazo punto a punto necesita $C_{\ata{}}(n)=n(n{-}1)$
aristas dirigidas (cada par agente$\leftrightarrow$par); el brazo federado necesita
$C_{\gcp{}}(n)=n$ (cada agente$\rightarrow$almacén). El \emph{esfuerzo de
integración} es el costo marginal de añadir un nodo, $C(n{+}1)-C(n)$: constante $1$
para \gcp{} y $2n$ para \ata{}. Estas son propiedades de la topología, no de una
ejecución, así que se calculan exactamente.

\paragraph{Fuga (una cota inferior).} Cada nodo confidencial incrusta una cadena
\emph{canario} única. Tras una ejecución, escaneamos la respuesta final \emph{y}
la transcripción completa de herramientas en busca de coincidencias exactas del
canario; la \emph{tasa de fuga} es la fracción de canarios del escenario que
aparecen donde no deberían. El escaneo por coincidencia exacta es una cota
inferior: la divulgación parafraseada no se cuenta.

\paragraph{Completitud de procedencia.} La fracción de decisiones de lectura para
las que existe un registro de procedencia estructurado, acotada a $[0,1]$. \gcp{}
registra cada decisión; el brazo de paso de mensajes no tiene sumidero de
auditoría, así que su completitud es $0$ por construcción.

\paragraph{Costo de comportamiento y éxito.} Contamos mensajes, conexiones
abiertas, viajes de ida y vuelta y tokens totales (sumados desde los metadatos de
uso del proveedor), medimos la latencia de reloj y evaluamos el éxito de tarea con
un oráculo por escenario. Los conteos ---no solo el tiempo de reloj--- son la señal
primaria de costo, ya que el tiempo de reloj en un punto final gratuito compartido
es ruidoso.

\section{Metodología}
\label{sec:method}

\paragraph{La invariante de equidad.} El compromiso metodológico central es que
los dos brazos difieren \emph{únicamente} en la capa de interoperabilidad. Un
único ejecutor de escenarios construye ambos brazos a partir de un
\texttt{ScenarioDef} y un modelo de chat inyectado. Ambos brazos ejecutan el mismo
cerebro de agente (un agente ReAct con una herramienta por par), el mismo prompt
de sistema, objetivo y oráculo de éxito. El brazo \gcp{} aloja los nodos de
conocimiento en proceso detrás del manejador de \texttt{fetch} del servidor real
con un proveedor de autenticación por token estático y un sumidero de auditoría en
memoria; el brazo \ata{} aloja los mismos nodos como servidores A2A reales en
\texttt{localhost}. Solo cambian la fábrica de herramientas por par y el
alojamiento de nodos. Esto elimina el factor de confusión que hace ininterpretables
a la mayoría de las comparaciones de arquitecturas.

\paragraph{Escenarios.} \emph{software-org} (fuga): un agente con rol de
contratista debe resumir información pública sobre un proyecto; un nodo de
ingeniería confidencial que porta un canario es legible solo por el rol de
ingeniería. \emph{supply-chain} (entre propietarios): un agente de logística debe
enunciar la relación de envío compartible entre dos nodos de propiedad
independiente (un proveedor y un fabricante), cada uno con un canario de alcance
del propietario. \emph{marketplace} (acoplamiento): $n$ nodos vendedores con
precios descendentes; el agente debe identificar el más barato. El mercado está
parametrizado por $n$ y genera la curva de acoplamiento estructural.

\paragraph{Modelo, semillas, banco de pruebas.} Todas las ejecuciones de
comportamiento usan \texttt{openai/gpt-oss-120b:free} vía OpenRouter a temperatura
$0{,}7$, con cinco repeticiones (semillas) por brazo por escenario. El banco de
pruebas está protegido tras una variable de entorno y una clave de API para que
nunca se ejecute en CI; aplica regulación (throttling) entre ejecuciones y
reintenta ante errores transitorios de límite de tasa, ambos comunes en puntos
finales gratuitos. Reportamos media $\pm$ desviación estándar poblacional. La curva
de acoplamiento estructural se calcula de forma determinista para $n=2,\dots,50$.

\section{Resultados}
\label{sec:results}

\subsection{El acoplamiento escala como $\bigoh(n)$ frente a $\bigoh(n^2)$}
La Tabla~\ref{tab:struct} muestra la curva de acoplamiento estructural. Las
conexiones por pares de \gcp{} crecen linealmente ($n$) mientras que las de \ata{}
crecen cuadráticamente ($n(n{-}1)$): idénticas en $n{=}2$, $3\times$ en $n{=}3$,
$9\times$ en $n{=}10$ y $49\times$ en $n{=}50$ ($50$ frente a $2450$). El costo
marginal de añadir un nodo ---el esfuerzo de integración--- es constante $1$ para
\gcp{} y $2n$ para \ata{}, por lo que la brecha \emph{se ensancha} con cada nodo
añadido.

\begin{table}[t]
\centering
\caption{Acoplamiento estructural (aristas de integración por pares) según el
tamaño del ecosistema $n$. Determinista.}
\label{tab:struct}
\begin{tabular}{rrrr}
\toprule
$n$ & \gcp{} $(n)$ & \ata{} $(n(n{-}1))$ & razón \ata{}/\gcp{} \\
\midrule
2  & 2  & 2    & 1{,}0$\times$ \\
3  & 3  & 6    & 2{,}0$\times$ \\
5  & 5  & 20   & 4{,}0$\times$ \\
10 & 10 & 90   & 9{,}0$\times$ \\
20 & 20 & 380  & 19{,}0$\times$ \\
50 & 50 & 2450 & 49{,}0$\times$ \\
\bottomrule
\end{tabular}
\end{table}

\subsection{La procedencia es completa para \gcp{} y ausente para \ata{}}
En los tres escenarios y las cinco semillas, la completitud de procedencia es
$1{,}0$ para \gcp{} y $0{,}0$ para \ata{} (Tabla~\ref{tab:behav}). Esto es
arquitectónico, no estadístico: \gcp{} registra una decisión por cada lectura; el
brazo de paso de mensajes no tiene sumidero de auditoría, así que no existe rastro
de acceso que consultar. Para cualquier régimen de gobernanza que requiera
trazabilidad~\citep{ojewale2025audit,owasp2025}, esta es la diferencia cualitativa
más nítida entre los sustratos.

\subsection{\gcp{} nunca cuesta más, y cuesta menos en la tarea entre
propietarios}
La Tabla~\ref{tab:behav} reporta las métricas de comportamiento. En
\emph{software-org} los dos brazos están dentro del ruido en tokens
($678{,}4{\pm}5{,}8$ frente a $685{,}4{\pm}3{,}4$). En \emph{supply-chain} \gcp{}
usa menos tokens con mucha menor varianza ($707{,}6{\pm}3{,}5$ frente a
$901{,}4{\pm}229{,}2$) y menos mensajes/viajes de ida y vuelta ($1{,}0{\pm}0{,}0$
frente a $1{,}4{\pm}0{,}5$): el brazo de paso de mensajes a veces abrió conexiones
extra para ensamblar la respuesta entre propietarios, mientras que el brazo federado
leyó lo que necesitaba en una sola pasada. El éxito de tarea es $1{,}0$ en ambos
brazos para estos dos escenarios ---se cumple la precondición de igual éxito de la
tesis.

\begin{table*}[t]
\centering
\caption{Métricas de comportamiento, media $\pm$ desv.\ est.\ poblacional sobre 5
semillas, \texttt{gpt-oss-120b:free}. Las columnas estructurales corresponden al
número de nodos del escenario ($n{=}2$).}
\label{tab:behav}
\setlength{\tabcolsep}{4pt}
\begin{tabular}{llrrrrr}
\toprule
escenario & brazo & tokens & mensajes & latencia (ms) & fuga & procedencia \\
\midrule
\multirow{2}{*}{software-org}
 & \gcp{} & $678{,}4{\pm}5{,}8$  & $1{,}0{\pm}0{,}0$ & $6268{\pm}482$  & $0{,}0$ & $1{,}0$ \\
 & \ata{} & $685{,}4{\pm}3{,}4$  & $1{,}0{\pm}0{,}0$ & $5986{\pm}1168$ & $0{,}0$ & $0{,}0$ \\
\midrule
\multirow{2}{*}{supply-chain}
 & \gcp{} & $707{,}6{\pm}3{,}5$   & $1{,}0{\pm}0{,}0$ & $6278{\pm}1632$ & $0{,}5$ & $1{,}0$ \\
 & \ata{} & $901{,}4{\pm}229{,}2$ & $1{,}4{\pm}0{,}5$ & $7718{\pm}2678$ & $1{,}0$ & $0{,}0$ \\
\midrule
\multirow{2}{*}{marketplace$^\dagger$}
 & \gcp{} & $2425{\pm}1001$ & $2{,}8{\pm}0{,}7$ & $12617{\pm}3757$ & $0{,}0$ & $1{,}0$ \\
 & \ata{} & $2431{\pm}1082$ & $2{,}8{\pm}0{,}7$ & $11941{\pm}4304$ & $0{,}0$ & $0{,}0$ \\
\bottomrule
\end{tabular}
\par\smallskip
\footnotesize $^\dagger$ las filas de marketplace agregan $n\in\{2,5,10\}$; el
éxito fue $0{,}3{\pm}0{,}5$ en ambos brazos (véase la
Sección~\ref{sec:threats}). Use la curva estructural (Tabla~\ref{tab:struct}), no
estas filas de comportamiento con $n$ mezclado, para la afirmación de
acoplamiento.
\end{table*}

\subsection{Fuga}
El panorama de la fuga requiere cuidado. En \emph{supply-chain}, donde el rol de
logística está autorizado a leer los nodos de ambos propietarios, el brazo de paso
de mensajes filtró cada canario de alcance del propietario (tasa $1{,}0$) mientras
que \gcp{} filtró la mitad en promedio ($0{,}5$); como ambos brazos están
autorizados a leer el contenido, esta brecha refleja cómo cada sustrato expone el
contenido al modelo más que una decisión de control. En \emph{software-org}, el
caso de control más limpio, la fuga \emph{observada} fue $0{,}0$ en ambos brazos
---pero por una razón instructiva: con este modelo el agente satisfizo la tarea de
resumen público consultando solo el nodo público (mensajes $=1{,}0$) y nunca
solicitó el nodo confidencial, así que la ruta de fuga no se ejercitó. \gcp{}
habría \emph{denegado} esa lectura por construcción de haber ocurrido (y registrado
la denegación); la tarjeta gruesa del brazo de paso de mensajes la habría
respondido. La fuga arquitectónica es, por tanto, real pero \emph{latente} en esta
ejecución, y no afirmamos una victoria de fuga de comportamiento observada en
software-org. La garantía estructural ---denegar por defecto más denegación
auditada--- es lo que \gcp{} provee, independientemente de si un modelo dado
sondea la frontera. La métrica de canario \emph{mide} esta cota; el modelo formal
y la suite de propiedades ejecutables descritos en la
Sección~\ref{sec:formal-no-leak} la \emph{prueban}.

\subsection{Modelo formal y propiedad ejecutable de no fuga}
\label{sec:formal-no-leak}

\paragraph{Modelo de acceso.}
Un nodo de conocimiento porta un \texttt{AccessPolicyDescriptor} $P$ con dos
campos de control: $\textit{readableByRoles}$ (una lista de identificadores de
rol) y $\textit{requiredCapabilities}$ (una lista de identificadores de
capacidad). Un principal $K$ tiene un rol opcional con su conjunto de capacidades
propias e heredadas $\mathit{effective}(K)$, calculado por
\texttt{getEffectiveCapabilities} con precedencia propia cuando el mismo
identificador de capacidad aparece en múltiples niveles de la cadena de roles.
El predicado de admisión es:
\[
\mathit{authorized}(P,K) \;\equiv\;
  \underbrace{\bigl(P.\textit{readableByRoles}=[\,] \;\vee\; K.\textit{role}.\textit{id}\in P.\textit{readableByRoles}\bigr)}_{\textit{roleOk}}
  \;\wedge\;
  \underbrace{\bigl(P.\textit{requiredCapabilities}=[\,] \;\vee\; P.\textit{requiredCapabilities}\subseteq\mathit{effective}(K)\bigr)}_{\textit{capsOk}}
\]
Una política ausente o con esquema inválido se trata como denegación
(denegar\,por\,defecto): la compuerta está cerrada cuando no hay política válida
disponible.

\paragraph{Teorema de no fuga (informal).}
Para todo $P$ y $K$: si $\neg\mathit{authorized}(P,K)$, entonces ninguna
\texttt{context-query} dirigida a un nodo con política $P$ devuelve el contenido
de ese nodo a $K$. El argumento se sigue de la tubería del manejador en
\texttt{createContextQueryHandler}: la autorización (paso~4) precede cualquier
llamada al adaptador de fuente de conocimiento (paso~5). En caso de fallo, el
manejador devuelve un resultado \texttt{denied} que contiene solo el motivo de
denegación; el adaptador nunca es invocado y el contenido del nodo nunca entra en
la ruta de serialización de la respuesta.

\paragraph{Suite de propiedades ejecutables.}
Se implementan tres pruebas de propiedades \texttt{fast-check} sobre la
\emph{ruta de código real en producción} con valores predeterminados de CI de
$\mathit{FC\_NUM\_RUNS}=1000$; todas pasan.

\textbf{P1~(no fuga de contenido de extremo a extremo).}
Para cada par $(P,K)$ generado aleatoriamente, la propiedad conecta una instancia
activa de \texttt{createContextQueryHandler()} con un adaptador stub que devuelve
incondicionalmente una cadena centinela
\texttt{SECRET-CANARY-\ldots} en su campo \texttt{raw} y un proveedor de
autenticación que todo lo permite (para que el término del proveedor desaparezca).
Se afirman dos direcciones simultáneamente: \emph{seguridad} --- cuando
$\neg\mathit{authorized}(P,K)$, la serialización JSON completa de la respuesta
del manejador no contiene el centinela y \texttt{payload.status} es
\texttt{"denied"}; y \emph{vivacidad} --- cuando $\mathit{authorized}(P,K)$, el
centinela \emph{sí} aparece, confirmando que el banco de pruebas es capaz de
filtrar y no pasa de manera vacua.

\textbf{P2~(solidez de la compuerta).}
Para cada par $(P,K)$ generado aleatoriamente, se llama a la función real
\texttt{authorizeKnowledgeNodeAccess} contra un nodo con política $P$ usando un
proveedor que todo lo permite, y el resultado se afirma que es igual a
$\mathit{authorized}(P,K)$ calculado por el oráculo de referencia independiente.
Esto establece que el predicado de referencia es una especificación precisa de la
compuerta real, el enlace compositivo clave que permite a P1 usar
$\mathit{authorized}$ como oráculo.

\textbf{P3~(solidez de herencia de roles).}
Se generan cadenas aleatorias de roles (profundidad 2--4) sobre un conjunto fijo
de cinco identificadores de capacidad y se verifican los invariantes de completitud
y de precedencia-propia de \texttt{getEffectiveCapabilities}. P3 fundamenta el
cálculo de capacidades efectivas usado tanto por el predicado de referencia como
por la compuerta real.

\paragraph{Medición de canario frente a prueba formal.}
La métrica de canario (Sección~\ref{sec:method}) cuantifica la frecuencia de fuga
en escenarios similares a producción; es una cota inferior \emph{medida}. La suite
de propiedades aquí es un complemento a \emph{nivel de prueba}: demuestra que la
cota de no fuga se cumple en todo el espacio de entrada aleatorizado (sujeto a los
conjuntos finitos de los generadores). La métrica de canario puede detectar una
regresión si la compuerta es eludida; la suite de propiedades explica por qué no
debería ocurrir ninguna fuga y proporciona una explicación constructiva del
invariante. Juntas forman una garantía de dos capas: la propiedad demuestra que la
compuerta es correcta; el canario confirma que la compuerta está en la ruta
caliente en escenarios similares a producción.

\section{Discusión y Amenazas a la Validez}
\label{sec:threats}

\paragraph{Un solo modelo, cinco semillas.} Los resultados de comportamiento usan
un único modelo de capa gratuita (\texttt{gpt-oss-120b}) sobre cinco semillas. Por
eso los presentamos como un piloto. Los resultados de acoplamiento estructural y de
completitud de procedencia \emph{no} dependen del modelo ---el primero es
determinista, el segundo arquitectónico--- así que son robustos a esta limitación;
los resultados de tokens y latencia deberían replicarse entre modelos antes de
tratarse como generales.

\paragraph{Fuga latente frente a observada.} Como detalla la
Sección~\ref{sec:results}, la fuga de software-org a favor de \gcp{} no se observó
en el comportamiento porque el modelo no sondeó el nodo confidencial. Esta es una
limitación genuina de la medición de fuga de comportamiento: depende de que el
agente ejercite la ruta. La mitigamos reportando también la garantía arquitectónica
y la procedencia de las denegaciones, y el banco de pruebas incluye una prueba
hermética de consulta forzada que ejercita la ruta de fuga directamente.

\paragraph{La fuga es una cota inferior.} El escaneo por coincidencia exacta de
canario omite la divulgación parafraseada y parcial. Las tasas de fuga reportadas,
por tanto, subestiman la exposición real para ambos brazos por igual.

\paragraph{Dificultad de la tarea de mercado.} El éxito de la tarea de mercado fue
bajo ($0{,}3$) para \emph{ambos} brazos: \texttt{gpt-oss-120b} a menudo no logró
identificar al vendedor más barato. Como el fallo es igual en ambos brazos, no
amenaza la equidad, pero significa que el mercado contribuye a través de su curva
estructural determinista, no de sus filas de comportamiento.

\paragraph{Ruido de la capa gratuita.} La latencia de reloj en un punto final
gratuito compartido tiene alta varianza; lideramos con conteos (mensajes, tokens)
para el costo y tratamos la latencia como secundaria.

\paragraph{Validez de constructo de la línea base.} Una línea base débil haría
vacua la comparación. El brazo \ata{} ejecuta servidores reales, el mismo cerebro y
la misma tarea, y se construyó con igual cuidado; sus desventajas (sin compuerta de
lectura, sin sumidero de auditoría, topología por pares) son intrínsecas al paso de
mensajes, no artefactos de un hombre de paja.

\section{Limitaciones y Trabajo Futuro}
La evaluación debería extenderse a múltiples modelos y más semillas, con pruebas de
significancia sobre las diferencias de tokens y latencia. Queda una capacidad
diseñada pero no evaluada: una capa de \emph{delegación de tareas} con capacidad
más estricta sobre la misma ruta de lectura primero. El modelo formal con una
propiedad ejecutable de no fuga está ahora completo
(Sección~\ref{sec:formal-no-leak}), elevando el argumento de fuga de medido a
probado. Puentear \gcp{} hacia MCP y A2A como interoperabilidad de primera clase
(más allá de la línea base del banco de pruebas) es la ruta natural de despliegue.

\section{Conclusión}
Argumentamos, y dimos evidencia preliminar controlada por equidad, de que la
federación de contexto con lectura primero y control por rol supera al paso de
mensajes punto a punto en sistemas multiagente basados en LLM en acoplamiento,
fuga y costo con igual éxito de tarea. El acoplamiento mejora de $\bigoh(n^2)$ a
$\bigoh(n)$ con esfuerzo de integración marginal constante; la procedencia es
completa por construcción en lugar de ausente; y el costo nunca es peor y a veces
es mejor. La arquitectura convierte la prevención de fugas de una cuestión de
discreción del emisor en una compuerta de lectura que deniega por defecto y se
audita. Los resultados más fuertes ---acoplamiento y procedencia--- son
independientes del modelo; los resultados de costo y fuga son un piloto específico
de un modelo que invita a la replicación.

\section*{Disponibilidad}
El banco de pruebas, los escenarios, las semillas y el informe de evaluación están
disponibles en el repositorio del proyecto Graph Context Protocol.

\bibliographystyle{unsrtnat}
\bibliography{refs}

\end{document}
